\section*{Appendix A. Formal Definitions of Evaluated Objectives}
\label{app:objectives}

This appendix gives the precise definition of each self-supervised objective evaluated in Section~\ref{sec:exp-objectives}. All objectives operate on the encoder output $\mathbf{Z}$ (Section~\ref{sec:method-encoder}) and use only information that survives missingness, namely graph structure, observed neighbor features, or observed training labels. We write $\mathcal{N}_k(i)$ for the $k$-hop neighborhood of node $i$, $V_o$ for the observable set, and $h_{(\cdot)}$ for a lightweight per-objective prediction head applied to $\mathbf{Z}$. Unless noted, $k$ matches the value used in training.

\paragraph{Feature-related objectives.} These predict per-node targets derived from observable neighbor features and are trained with a mean-squared error.
The \emph{centroid} objective predicts the mean of a node's observable neighbors' features (for nodes with a non-empty observable neighborhood),
\begin{equation}
    \mathbf{t}^{\text{cen}}_i = \frac{1}{|\mathcal{N}_k(i)\cap V_o|}\!\!\sum_{j\in\mathcal{N}_k(i)\cap V_o}\!\!\mathbf{x}_j,
    \qquad
    \mathcal{L}_{\text{cen}} = \frac{1}{n}\sum_i \big\|h_{\text{cen}}(\mathbf{z}_i)-\mathbf{t}^{\text{cen}}_i\big\|_2^2 .
\end{equation}
The \emph{stats} objective is identical but predicts the per-dimension standard deviation of the same neighborhood, $\mathbf{t}^{\text{std}}_i = \mathrm{std}_{j\in\mathcal{N}_k(i)\cap V_o}\,\mathbf{x}_j$.
The \emph{recon} objective is a masked-feature self-reconstruction: at each step half of the observable nodes are held out, their features zeroed, and their true features regressed from the re-encoded representations,
\begin{equation}
    \mathcal{L}_{\text{rec}} = \frac{1}{|H|}\sum_{i\in H}\big\|h_{\text{rec}}(\mathbf{z}_i)-\mathbf{x}_i\big\|_2^2 ,
\end{equation}
where $H\subset V_o$ is the freshly sampled held-out set.

\paragraph{Structural objectives.} These predict topological targets computed purely from the graph and are defined even when no features are observed.
The \emph{path} objective is a stochastic multi-hop link prediction. A fresh batch of random walks of length $\ell=2$ yields positive pairs $(s,e)$ with $e$ reachable within $\ell$ hops of $s$; negatives $(s,e^-)$ draw $e^-$ from the set of nodes at least $3$ hops from $s$. Since $\ell<3$, the positive and negative sets do not overlap. With a scoring head $h_{\text{path}}$ on the elementwise product of endpoints,
\begin{equation}
    \mathcal{L}_{\text{path}} = \tfrac{1}{2}\Big[
        \mathrm{BCE}\big(h_{\text{path}}(\mathbf{z}_s\!\odot\!\mathbf{z}_e),\,1\big)
        + \mathrm{BCE}\big(h_{\text{path}}(\mathbf{z}_s\!\odot\!\mathbf{z}_{e^-}),\,0\big)
    \Big].
\end{equation}
The \emph{triplet} objective ranks distances in embedding space: for an anchor $a$, a near node $p$ (within $\ell$ hops) and a far node $q$ ($\geq 3$ hops), it enforces a margin,
\begin{equation}
    \mathcal{L}_{\text{tri}} = \big[\,\|\mathbf{z}_a-\mathbf{z}_p\|_2^2 - \|\mathbf{z}_a-\mathbf{z}_q\|_2^2 + 1\,\big]_+ .
\end{equation}
The \emph{anchor} objective regresses each node's normalized hop-distances to $K$ sampled landmark nodes, $\mathbf{d}_i\in[0,1]^K$, computed by multi-source breadth-first search; \emph{anchorcls} is its classification variant, predicting the discrete hop-distance bucket to each landmark with a cross-entropy loss.

\paragraph{Label-related objective.} The \emph{hist} objective is semi-supervised. Its target is the class distribution of a node's observable $k$-hop neighbors, using only training labels,
\begin{equation}
    \mathbf{p}_i = \mathrm{normalize}\Big(\sum_{j\in\mathcal{N}_k(i)\cap V_o} w_{ij}\,\mathbf{e}_{y_j}\Big),
\end{equation}
where $\mathbf{e}_{y_j}$ is the one-hot label of neighbor $j$ and $w_{ij}$ a proximity weight. The prediction is trained with a KL divergence over observable nodes that have at least one observable neighbor,
\begin{equation}
    \mathcal{L}_{\text{hist}} = \frac{1}{|V_o|}\sum_{i\in V_o} \mathrm{KL}\big(\mathbf{p}_i \,\|\, \mathrm{softmax}(h_{\text{hist}}(\mathbf{z}_i))\big).
\end{equation}

\paragraph{Contrastive objectives.} These do not predict a fixed target but encourage agreement between the two views $\mathbf{Z}^{(1)},\mathbf{Z}^{(2)}$ (Section~\ref{sec:method-encoder}). The default is the Barlow-Twins redundancy-reduction loss of Equation~\eqref{eq:con}. We additionally evaluate \emph{InfoNCE}, which maximizes agreement between the two views of the same node against other nodes in the batch, and \emph{prototype-InfoNCE}, which contrasts nodes against per-class prototypes formed from observable labels.






\section*{Appendix B. Hyperparameters and Design Choices}
\label{sec:setup-hparams}

Table~\ref{tab:hparams} lists the full NESS configuration used in all experiments, grouped by the two feature regimes: the large dense-embedding graph (ogbn-arxiv) and the small sparse binary bag-of-words graphs (Cora, CiteSeer, AMAP, AMAC). The architecture, self-supervised objectives, and loss weights are shared across both regimes; the settings that differ are the learning rate and the use of cluster-based batching, both dictated by graph scale rather than by tuning. These values were selected once on ogbn-arxiv (Section~\ref{sec:exp-design}) and held fixed thereafter.

\begin{table}[t]
\centering
\caption{NESS hyperparameters and design choices. The architecture and objectives are shared; only the scale-dependent settings (learning rate, batching) differ between regimes.}
\label{tab:hparams}
\small
\begin{tabular}{@{}lcc@{}}
\toprule
\textbf{Setting} & \textbf{ogbn-arxiv} & \textbf{Binary BoW graphs} \\
 & (dense) & (Cora/CiteSeer/AMAP/AMAC) \\
\midrule
\multicolumn{3}{@{}l}{\emph{\textbf{Architecture}}} \\
Prefill                & Feature Propagation & Feature Propagation \\
Prefill iterations $T$ & 40 & 40 \\
Encoder                & GraphSAGE & GraphSAGE \\
Layers $L$             & 3 & 3 \\
Hidden width $h$       & 128 & 128 \\
Dropout                & 0.3 & 0.3 \\
Encoder / decoder channels & 256 / 64 & 256 / 64 \\
\midrule
\multicolumn{3}{@{}l}{\emph{\textbf{Views and objectives}}} \\
View 1 / View 2        & edge-mask / edge-mask & edge-mask / edge-mask \\
Edge-mask rate $p$     & 0.7 & 0.7 \\
View fusion            & average & average \\
Contrastive            & Barlow-Twins ($\lambda=0.005$) & Barlow-Twins ($\lambda=0.005$) \\
SSL objectives         & hist + path & hist + path \\
SSL neighborhood $k$   & 2 & 2 \\
Path walk length $\ell$ & 2 & 2 \\
\midrule
\multicolumn{3}{@{}l}{\emph{\textbf{Loss weights}}} \\
$w_{\text{cls}}$       & 3.0 & 3.0 \\
$w_{\text{edge}}$      & 1.0 & 1.0 \\
$w_{\text{con}}$       & 0.5 & 0.5 \\
$w_{\text{ssl}}$ (hist, path) & 1.0, 1.0 & 1.0, 1.0 \\
\midrule
\multicolumn{3}{@{}l}{\emph{\textbf{Optimization}}} \\
Optimizer              & Adam & Adam \\
Learning rate          & $10^{-3}$ & $10^{-2}$ ($10^{-3}$ for AMAC) \\
Weight decay           & $5\times10^{-5}$ & $5\times10^{-5}$ \\
Epochs                 & 800 & 800 \\
Early stopping         & val macro-F1, patience 20 & val macro-F1, patience 20 \\
Batching               & cluster-based (50 parts) & full-batch \\
\bottomrule
\end{tabular}
\end{table}
